\section{Pruebas focales de alertas persistentes}\label{sec:pruebas-alertas}

La verificación de alertas se organiza por fronteras de responsabilidad, con los resultados del corte funcional del 19 de septiembre de 2026 registrados en \rutarepo{docs/verification-report.md}; la familia propia de las audiencias de recursos tiene su verificación en la \cref{sec:pruebas-alertas-audiencias-recursos}. Ocho pruebas de temporización y 16 de aplicación comprobaron ventanas, preferencias y contratos; tres pruebas de vigencia separaron revisión humana y recálculo ordinario, y seis pruebas HTTP comprobaron el transporte. El cliente aprobó 26 pruebas, seguidas de una prueba específica que detectó la aceptación de un cursor futuro y aprobó después de corregir ese comportamiento. Nueve recorridos de navegador con transporte HTTP controlado aprobaron en 18.891 segundos e incluyeron escritorio, móvil, respuesta incierta, conflicto, apertura exacta, retorno, continuación y revocación. La inspección de dos capturas de bandeja no detectó recortes ni solapamientos. Estas comprobaciones no se presentan como aceptación de servicios reales.

La primera comprobación de PostgreSQL real utilizó bases desechables y un único proceso de verificación. Veinte casos aprobaron y dos reprodujeron defectos en la continuidad de episodios y la reactivación de preferencias. Los casos aprobados comprendieron idempotencia, aislamiento antes del límite, reinicio, lectura independiente de atención, origen histórico corrupto, reversión por fallo de auditoría, fuente modificada antes de la reevaluación, intentos de correo, inventario y privilegios mínimos. El rechazo de alteraciones del catálogo incluyó restricciones, índices, funciones, disparadores y permisos públicos o heredados. La campaña no permite declarar aprobadas las dos conductas que aún fallaban en esa ejecución.

La corrección conservó el cierre del episodio al registrar atención, incluso si vuelve a pendiente antes del siguiente escaneo, y permitió reactivar planes que nunca habían emitido un aviso al recuperar sus canales. Ambas pruebas aprobaron en la repetición focal. Una tercera prueba reprodujo la pérdida equivalente de episodio cuando una aceptación humana es seguida por otro cambio de fuente antes del escaneo. Tras conservar esa transición y verificar la revisión histórica al resolver el aviso anterior, los tres casos aprobaron juntos en 26.463 segundos. Se comprobó la nueva ocurrencia activa, la resolución del aviso precedente y la ausencia de duplicados después de reabrir el almacén. Esta repetición no añade de nuevo los dos casos anteriores a un total global.

El adaptador de correo aprobó seis pruebas contra un servidor local de ensayo: solicitud e idempotencia, aceptación explícita, incertidumbre, rechazo y reintentos. No se contactó a un destinatario externo. Ocho pruebas previas de configuración, límites de ciclo, parada y supervisión comprobaron la composición mediante puertos antes de conectar PostgreSQL. Posteriormente aprobaron dos pruebas de arranque del servidor, incluidas la imposibilidad de iniciar consumidores cuando falla la apertura del puerto de red y la parada conjunta con liberación de adaptadores fuera de Tokio. Cuatro pruebas de ayuda y límites de argumentos de la CLI también aprobaron. Cada grupo conserva su alcance y ejecución independiente.

La aceptación de navegador con servicios reales aprobó tres escenarios: bandeja de alertas a 1440 y 390 píxeles y permisos con revocación. El comando completo tardó 233.822 segundos, incluida preparación; Playwright informó 26.7 segundos. Los recorridos comprobaron generación por el consumidor, origen exacto, persistencia de lectura, preferencias y transporte de correo deshabilitado. Se inspeccionaron las dos capturas de la bandeja, sin desbordamientos ni solapamientos. Esta comprobación utilizó el servidor compilado antes de una corrección posterior de restauración, cuyo resultado se registra por separado.

La aceptación final mediante la API aprobó en 273.199 segundos, con las fuentes sin cambios. Comprobó generación real, bandeja propia, preferencias, lectura idempotente y denegaciones. El respaldo comparó exactamente las ocho tablas de alertas antes y después de restaurar con el servidor detenido. Al reabrirlo, las consultas conservaron preferencias, aviso leído, destinatario, origen e historia de audiencia, sin duplicar la ocurrencia. Dos pruebas focales adicionales comprobaron la validación SQL con espacio de búsqueda vacío y la ausencia de escrituras cuando el escáner no encuentra actividades. Este corte no acredita recepción externa de Resend ni cierre de la regresión completa. La generación de un aviso tampoco demuestra validez jurídica del acto, presentación oportuna, identidad jurídica acreditada ni usabilidad con personas. La CA interna y la TSA local mantienen sus límites experimentales. El informe \rutarepo{docs/verification-report.md} conserva las ejecuciones y sus resultados; el \cref{sec:matriz-alertas} relaciona las fuentes correspondientes.

Una comprobación posterior de reinicio detectó un plan de audiencia incompatible con el inventario: al cambiar una fecha próxima, el planificador generaba un aviso de cambio de vencimiento reservado a plazos. La reproducción focal produjo dos fallos y un control de plazo aprobado; con PostgreSQL aislado se observaron los mismos dos fallos y un control aprobado en 14.55 segundos. El caso persistido fallaba al activar el aviso o al reabrir el almacén, después de una parada normal del servidor. La corrección restringe el episodio de cambio a plazos; una audiencia reprogramada conserva únicamente anticipaciones con fecha y origen vigentes, conforme a \rutarepo{docs/alerts-api.md}.

Después de la corrección aprobaron tres casos unitarios y tres PostgreSQL, estos últimos en 14.96 segundos. Son seis casos distintos, correspondientes a la misma reproducción, y no pruebas adicionales por repetirla. Incluyen planes pendientes, avisos activados, reapertura sin reescritura de filas y conservación del aviso de cambio de vencimiento de un plazo. Esa comprobación focal no repitió el reinicio del servidor compuesto. La aceptación HTTP de informes del 2 de octubre de 2026 registra reinicios con el inventario poblado mediante \texttt{SIGTERM} y \texttt{SIGINT}, que conservaron el histórico exacto; los ensayos de arranque previos que terminaron antes de generar informes, como el que detectó este fallo, no se cuentan como esa aceptación (\cref{sec:pruebas-informes}). La prevención conserva las validaciones de activación e inventario; no borra ni reclasifica registros ya inconsistentes para permitir el arranque. Estos resultados no se suman a los conteos históricos ni acreditan una reparación de datos existentes.

\subsection{Verificación focal de alertas de audiencias de recursos}\label{sec:pruebas-alertas-audiencias-recursos}

La familia propia descrita en la \cref{sec:impl-alertas-audiencias-recursos} se verificó el 4 de octubre de 2026. El contrato de aplicación incorporó cuatro casos nuevos escritos antes de la variante de sujeto. Verifican identidad propia frente a UUID compartidos, origen inicial y motivo de proximidad, límites temporales exactos y conservación del recibo de lectura. Los cuatro y las dieciséis regresiones de la suite aprobaron. Tres comprobaciones HTTP nuevas y seis anteriores aprobaron proyección del recurso padre, captura e instante, lectura y rechazo cerrado de una revisión o motivo impropios.

El cliente verificó seis casos nuevos y veintiséis anteriores. Una fixture de preferencias modificaba los valores iniciales sin recibo y fue rechazada correctamente; al representarla como revisión persistida con recibo y fecha, ese caso aprobó sin modificar el validador. Los treinta y dos casos distintos quedan acreditados por la campaña parcial y su corrección focal, no como dos regresiones completas.

Cinco recorridos nuevos de navegador con HTTP controlado aprobaron en 16.9 segundos y once regresiones anteriores en 27.8 segundos, con un trabajador. Comprueban apertura histórica a 1440 y 390 píxeles, permiso retirado, huella discordante y descarte de una respuesta tardía tras cambiar filtros. La repetición de dos positivos en 10.2 segundos renovó únicamente las capturas visuales y no añade escenarios distintos. Se conserva la separación entre pruebas de cliente, persistencia real y aceptación integrada del producto; estos resultados no acreditan entrega externa de correo, una restauración completa, cobertura global ni despliegue.

La verificación de persistencia empleó PostgreSQL 16.15 desechable. Quince casos distintos comprobaron columnas y permisos, parentesco exacto, preferencias, destinatarios, reinicio durante la paginación y ausencia de duplicados. Catorce aprobaron en la primera ejecución compilable; la simulación restante fue detenida por una restricción SQL antes de alcanzar el adaptador. Tras conservar una huella coherente con los bytes deliberadamente alterados, el caso comprobó el rechazo del origen reconstruido sin escrituras. Las veintiocho regresiones previas de alertas aprobaron en 200.15 segundos, sin repetir los quince casos nuevos. Tres pruebas adicionales conservaron la representación de las familias anteriores y rechazaron padres malformados. Las pruebas de reapertura e inventario no sustituyen la aceptación completa de respaldo y restauración, que conservó la misma ocurrencia y su lectura después de reiniciar el generador (\cref{sec:pruebas-audiencias-recursos}).
